\documentclass[11pt]{article}
\usepackage[margin=1in]{geometry}
\usepackage{amsmath,amssymb}
\usepackage{booktabs}
\usepackage{graphicx}
\usepackage[hidelinks]{hyperref}
\usepackage{times}

\title{\textbf{The Kuramoto Engine}\\[4pt]
\large Physics and Mathematics of the Cardiac Synchronization Index\\[2pt]
\normalsize Technical reference for the NEUROCARDIA application\\[2pt]
\small Version 0.2 --- Revision}
\author{Americo Simoes \\ \texttt{amexsimoes@gmail.com}}
\date{October 14, 2026}

\begin{document}
\maketitle

\begin{center}
\small
Copyright \textcopyright{} 2026 Americo Simoes. All rights reserved.\\
This document is released under the Creative Commons Attribution 4.0
International License (CC-BY-4.0). The software it describes is released
under the GNU General Public License, version 3 or later (GPL-3.0-or-later).
\end{center}

\vspace{1em}
\noindent\fbox{\parbox{\dimexpr\linewidth-2\fboxsep-2\fboxrule}{%
\textbf{Revision history.} This is revision 0.2 of \emph{The Kuramoto
Engine}. Revision 0.1 (October 14, 2026) specified the model, the
numerical scheme, and the parameter choices, and stated $K = 2$ as an
empirical choice. Revision 0.2 adds the result of the coupling-strength
sweep: $K = 2$ is not merely the value that produces a diagnostic
residual, it is the only value in the tested range where the residual
carries information at all. The revision also notes that the trajectory
of the order parameter is not robust to the frequency distribution at
any value of $K$ where the metric has dynamic range. The mathematical
treatment of the model is unchanged.}}

\vspace{1em}

\begin{abstract}
This document specifies the mathematical model, the numerical methods,
and the parameter choices underlying the Cardiac Synchronization Index
(CSI) computed by the NEUROCARDIA application. The CSI is the order
parameter of a Kuramoto coupled-oscillator network driven by the RR
interval series of a single-lead ECG recording. We present the model,
the discretization scheme, the treatment of the natural frequency
distribution, the phase-reset mechanism, and the derivation of the four
early warning signal metrics that accompany the CSI. Section~4.3 reports
the coupling-strength sweep and the narrow window in which the metric is
diagnostic. Section~4.2 discusses the sensitivity of the trajectory to
the frequency distribution. The document is self-contained: it assumes
no prior familiarity with the Kuramoto model.
\end{abstract}

\tableofcontents

\section{Introduction}

The Kuramoto model is the canonical mathematical description of
synchronization in a population of coupled phase oscillators. It was
introduced by Yoshiki Kuramoto in 1975~\cite{kuramoto1975} as a minimal
model of collective behavior, and it has since been applied across
physics, biology, and neuroscience~\cite{acebron2005,strogatz2000}. The
appeal of the model is that it captures the essential phenomenon ---
oscillators settling into or falling out of phase lock --- without
requiring detailed specification of the underlying dynamics.

This document describes one specific application of the Kuramoto model:
the analysis of cardiac rhythm from the RR interval series of a
single-lead ECG. The sinoatrial node is a population of coupled
pacemaker cells whose collective firing produces the cardiac rhythm.
Treating the RR tachogram as a phase signal and applying the Kuramoto
model to it yields a time-varying order parameter that we call the
Cardiac Synchronization Index (CSI). The CSI is a measure of how
regularly the heart is beating, expressed in the phase domain rather
than the time or frequency domain.

The document is organized as follows. Section~\ref{sec:model} presents
the Kuramoto model and the order parameter. Section~\ref{sec:rrphase}
describes the RR-to-phase mapping that connects the model to the ECG
signal. Section~\ref{sec:network} specifies the network parameters,
including the frequency distribution and the coupling strength.
Section~\ref{sec:integration} presents the numerical integration scheme.
Section~\ref{sec:ews} derives the early warning signal metrics.
Section~\ref{sec:properties} discusses mathematical properties of the
model. Section~\ref{sec:limitations} identifies the model's limitations.

\section{The Kuramoto Model}\label{sec:model}

\subsection{The phase oscillator}

Consider a system of $N$ coupled oscillators. Each oscillator $i$ is
described by a single scalar, its phase $\theta_i \in \mathbb{R}/2\pi\mathbb{Z}$.
The phase advances monotonically in time; the rate of advance is the
oscillator's instantaneous frequency. When the oscillator is isolated,
the rate is a constant called the natural frequency $\omega_i$.

The choice to describe an oscillator by its phase rather than by a full
state vector is a deliberate simplification. It is justified when the
oscillator's amplitude is either fixed or irrelevant to the question
being asked, and when the only observable of interest is timing. For
coupled pacemaker cells driving a cardiac rhythm, both conditions hold:
the observable that matters is when the next beat occurs, not the
detailed ionic dynamics that produced it.

\subsection{The coupling}

When two oscillators are coupled, each one's instantaneous frequency is
perturbed by the phase of the other. A natural form for this
perturbation is the sine of the phase difference. This choice has three
properties that make it the canonical one:

\begin{enumerate}
\item \emph{Antisymmetry.} If oscillator $j$ advances the phase of
oscillator $i$ by $\sin(\theta_j - \theta_i)$, then $i$ advances $j$ by
$\sin(\theta_i - \theta_j) = -\sin(\theta_j - \theta_i)$, the exact
opposite. The two oscillators influence each other equally and
oppositely. No energy is created or destroyed by the coupling.
\item \emph{Smoothness.} The sine is infinitely differentiable and
bounded. It cannot produce abrupt jumps in frequency.
\item \emph{$2\pi$-periodicity.} The coupling depends only on the phase
difference modulo $2\pi$. Two oscillators at phase difference $0.1$ and
at phase difference $0.1 + 2\pi$ are coupled identically.
\end{enumerate}

More general coupling functions are possible and are studied in the
literature. The sine coupling is the simplest that satisfies the
requirements above, and it is the form used here.

\subsection{The governing equation}

The Kuramoto model is the system of $N$ coupled differential equations
\begin{equation}
\frac{d\theta_i}{dt} = \omega_i + \frac{K}{N}\sum_{j=1}^{N}
\sin(\theta_j - \theta_i), \quad i = 1, \ldots, N,
\label{eq:kuramoto}
\end{equation}
where $K \geq 0$ is the coupling strength. The factor $1/N$ normalizes
the sum so that the total influence on oscillator $i$ does not grow
without bound as $N$ increases. With this normalization, the coupling
term is the average of the sine of the phase differences between $i$
and every other oscillator.

The equation has a clear physical reading. The phase of oscillator $i$
advances at its natural frequency $\omega_i$, plus a correction from
the coupling. If $i$ is lagging behind the population on average, the
sum of sines is positive and $i$ speeds up. If $i$ is leading, the sum
is negative and $i$ slows down. The coupling pulls oscillators toward
each other.

\subsection{The order parameter}

The degree to which the population is synchronized is quantified by the
order parameter
\begin{equation}
r = \left| \frac{1}{N} \sum_{j=1}^{N} e^{i\theta_j} \right|,
\label{eq:orderparam}
\end{equation}
where $e^{i\theta_j}$ is the point on the unit circle corresponding to
phase $\theta_j$. The order parameter has a direct geometric
interpretation: it is the magnitude of the mean of $N$ unit vectors,
each pointing at the phase of one oscillator.

If all oscillators have the same phase, the mean of their unit vectors
is a unit vector, and $r = 1$. If the phases are uniformly distributed
around the circle, the unit vectors sum to zero, and $r = 0$. For a
finite population of $N$ oscillators with uniformly random phases, the
expected value of $r$ is $1/\sqrt{N}$.

The order parameter is bounded in $[0,1]$ and is invariant to a global
phase shift: adding a constant to every $\theta_j$ does not change $r$.

\subsection{The critical coupling}

For a given distribution of natural frequencies $g(\omega)$, the
Kuramoto model exhibits a phase transition at a critical coupling
strength $K_c$. Below $K_c$, the oscillators drift incoherently and
$r \to 0$ in the thermodynamic limit. Above $K_c$, a subset of
oscillators locks to a common frequency and $r$ rises toward a positive
value.

For a symmetric, unimodal distribution $g(\omega)$, the critical
coupling has the closed form
\begin{equation}
K_c = \frac{2}{\pi g(0)},
\label{eq:critical}
\end{equation}
where $g(0)$ is the density of the frequency distribution at its
center. Distributions with a sharper peak at zero have smaller $K_c$;
distributions that are flatter at zero have larger $K_c$.

This closed form does not apply to the frequency distribution used here
(see Section~\ref{sec:network}), but the qualitative behavior holds:
there is a threshold below which the network is incoherent and above
which it synchronizes. The precise value of the threshold depends on
the distribution.

\section{RR-to-Phase Mapping}\label{sec:rrphase}

The Kuramoto model describes the evolution of phases. The ECG does not
measure phases directly. We construct a phase signal from the RR
interval series as follows.

Let $\{R_k\}_{k=1}^{M}$ be the intervals between successive detected
R-peaks, in seconds. Let $\overline{R}$ be their mean. Define the
cumulative phase at beat $k$ as
\begin{equation}
\phi_k = 2\pi \sum_{j=1}^{k} \frac{\overline{R}}{R_j} \pmod{2\pi}.
\label{eq:rrphase}
\end{equation}

Each beat advances the cumulative phase by $2\pi \overline{R}/R_j$. A
beat at the mean rate ($R_j = \overline{R}$) advances the phase by
exactly $2\pi$, returning to the same point on the circle. A shorter
interval advances the phase by more than $2\pi$; a longer interval
advances it by less.

The choice of $\overline{R}$ as the reference is deliberate. It makes
the phase trajectory scale-invariant: a subject whose mean RR is 0.6
seconds (100 BPM) and a subject whose mean RR is 1.2 seconds (50 BPM)
produce comparable phase trajectories, both advancing by $2\pi$ per
mean beat. The CSI measures synchronization, not rate. Rate is reported
separately as heart rate in beats per minute.

The cumulative phase is defined modulo $2\pi$. The natural choice is to
wrap at the boundary of the unit circle, mapping every phase into
$[0, 2\pi)$.

\section{Network Construction}\label{sec:network}

\subsection{Number of oscillators}

The network has $N = 16$ oscillators. This is a modeling choice. The
choice is not derived from cardiac physiology; it matches the dimension
used in the reference implementation from which the parameterization
was inherited. A principled choice --- for example, derived from the
number of identifiable frequency components in the RR spectrum, or from
the number of pacemaker cell clusters in the sinoatrial node --- would
be preferable, and the sensitivity of the CSI to this choice is a
subject for future work.

The incoherent baseline for $N = 16$ is $1/\sqrt{16} = 0.25$. Any CSI
value meaningfully above 0.25 indicates partial synchronization.

\subsection{Natural frequency distribution}

The natural frequencies $\omega_i$ are drawn from the non-trivial zeros
$\gamma_i$ of the Riemann zeta function. The first 24 such zeros are
tabulated in Table~\ref{tab:zeros}. For a network of $N = 16$
oscillators, the first 16 zeros are used.

\begin{table}[h]
\centering
\small
\begin{tabular}{llllll}
\toprule
14.134725 & 21.022040 & 25.010858 & 30.424876 & 32.935062 & 37.586178 \\
40.918719 & 48.005151 & 49.773832 & 52.970321 & 56.446248 & 59.347044 \\
60.831779 & 65.112544 & 67.079811 & 69.546402 & 72.067158 & 75.704691 \\
77.144840 & 79.337375 & 82.910381 & 84.735493 & 86.970000 & 87.425275 \\
\bottomrule
\end{tabular}
\caption{The first 24 non-trivial Riemann zeros. Used as the raw
frequencies of the network before normalization.}
\label{tab:zeros}
\end{table}

The raw frequencies are then normalized:
\begin{equation}
\omega_i = \frac{\gamma_i - \bar\gamma}{\sigma_\gamma},
\label{eq:normalization}
\end{equation}
where $\bar\gamma$ and $\sigma_\gamma$ are the mean and standard
deviation of the first 16 Riemann zeros.

Two points deserve emphasis.

\begin{enumerate}
\item The Riemann zeros are not claimed to have cardiac significance.
They are a deterministic, fixed sequence with a mathematically
determined distribution. Any other deterministic sequence with
comparable spread would serve the same role. The choice reflects the
reference implementation rather than a physiological argument.

\item The normalization is required for numerical tractability.
Section~\ref{sec:integration} discusses the integration scheme. The
unnormalized frequencies span a range of order $10^2$, which would
require an impractical step size for stable Euler integration. The
normalized frequencies span a range of order 1, which the Euler scheme
resolves at a tractable step.
\end{enumerate}

\paragraph{Sensitivity to the frequency distribution.}
The trajectory of the order parameter depends on the choice of frequency
distribution. A sensitivity analysis across five distributions
(Riemann zeros, uniform random, Gaussian, physiological scale, and
linear spacing) shows that the mean Pearson correlation between the CSI
trajectory under the Riemann-zero configuration and under any
alternative configuration is at most $0.17$ at $K = 2$, and reaches the
pre-registered threshold of $0.80$ only for the linear-spacing
distribution at $K = 30$, where the network is saturated. The absolute
value of the CSI is therefore not robust to the frequency distribution
at any coupling strength where the metric has dynamic range. The
residual diagnostic (Section~\ref{sec:properties}) is partially robust,
but not to the extent that revision 0.1 of the companion paper claimed.
We defer the full analysis to the companion documents and note here only
that the choice of frequency distribution is part of the definition of
the CSI, not an incidental parameter.

\subsection{Coupling strength}

The coupling strength is set to $K = 2$. This value is chosen
empirically rather than from the critical coupling formula, and it
matters: the behavior of the metric depends critically on this choice.
At $K = 2$ the network partially synchronizes after each R-peak reset
without saturating, so the order parameter trajectory retains dynamic
range. At $K \geq 5$ the network locks completely after each reset and
the order parameter sits near $1.0$ for all input rhythms, destroying
sensitivity. At $K \leq 1$ the network is under-coupled and the residual
diagnostic information inverts in sign.

The useful range is narrow and centered on $K = 2$. The companion
sensitivity analysis and residual papers report the full coupling sweep.
The quantitative summary is:

\begin{itemize}
\item \textbf{$K \leq 1$ (under-coupled).} The residual separates sinus
from AF with opposite sign: AUC $0.17$ at $K = 0.5$ and $0.21$ at
$K = 1.0$ (sign-inverted accuracies of $83\%$ and $79\%$). The network
is too weakly coupled to reflect the phase perturbations from the RR
series.
\item \textbf{$K = 2$ (intermediate).} The residual separates sinus from
AF with AUC $0.89$, $\rho = +0.68$, $p = 0.002$ on the pilot corpus.
This is the useful regime.
\item \textbf{$K \geq 3$ (saturated).} The residual AUC drops to
$0.59$--$0.72$ and the metric becomes indistinguishable from noise.
The network locks completely after each reset and stays locked until the
next beat, so the order parameter does not respond to the RR series.
\end{itemize}

The choice of $K$ is part of the metric's definition, not a physical
postulate about the sinoatrial node.

\subsection{Coupling topology}

The default coupling is all-to-all with weight $1/N$. This is the
classical Kuramoto configuration. Every oscillator couples to every
other with equal weight, and the topology reflects the assumption that
the sinoatrial pacemaker cells are densely interconnected.

Three alternative topologies are available in the implementation:

\begin{itemize}
\item \textbf{Nearest-neighbor.} Each oscillator couples to its two
immediate neighbors. Produces traveling waves.
\item \textbf{Small-world.} A ring with a small number of random
long-range links. Balances local clustering with global connectivity.
\item \textbf{Star.} One hub oscillator couples to all leaves.
Concentrates synchronization through a single node.
\end{itemize}

The topology is a modeling choice. The default is all-to-all.

\section{Numerical Integration}\label{sec:integration}

\subsection{The Euler scheme}

The Kuramoto equation~\eqref{eq:kuramoto} is a system of first-order
ordinary differential equations with no closed-form solution for
$N > 2$. We integrate numerically using the explicit Euler method:
\begin{equation}
\theta_i(t + \Delta t) = \theta_i(t) + \Delta t \cdot \left[
\omega_i + \frac{K}{N} \sum_{j=1}^{N} \sin(\theta_j(t) - \theta_i(t))
\right],
\label{eq:euler}
\end{equation}
where $\Delta t$ is the integration step. The phases are updated
simultaneously: the right-hand side is evaluated at the current time
$t$ for all $i$, then all updates are applied.

The choice of the Euler method is deliberate. Higher-order methods
(Runge-Kutta, for instance) are more accurate per step, but the Euler
method is stable for the step sizes used here and its behavior is easy
to reason about. The accuracy of the trajectory is not the objective;
the value of the order parameter is, and that quantity is robust to
moderate integration error.

\subsection{Step size}

The step size must be small enough that the phase advance per step is
much less than $2\pi$. Otherwise the sine evaluation aliases: the phase
advances past one full rotation during a single step, and the update
effectively samples a random point on the circle.

With normalized frequencies of order 1 and coupling $K = 2$, the
maximum instantaneous phase velocity is of order 2 rad/s. A step of
$\Delta t = 0.05$ s advances the phase by up to 0.5 radians per step,
well under the aliasing threshold.

The step size is fixed at $0.05$ seconds. For a recording with mean RR
interval $\overline{R}$, the number of integration steps per beat is
\begin{equation}
n_{\text{steps}} = \left\lfloor \frac{\overline{R}}{\Delta t} \right\rfloor,
\label{eq:steps}
\end{equation}
clamped to the range $[1, 200]$. For a mean RR of $0.83$ seconds (72
BPM), this gives 16 steps per beat.

\subsection{Phase reset}

At each detected R-peak, the phase of a designated heart oscillator
(index 0) is reset to the cumulative phase $\phi_k$ defined in
equation~\eqref{eq:rrphase}. The other $N-1$ oscillators are not reset;
they continue their evolution from the current state.

The reset is the mechanism by which the network is coupled to the
cardiac signal. Without it, the network would evolve as a free-running
oscillator system with no relationship to the subject's rhythm. With
it, each beat injects a phase impulse into the network, and the
network's response --- how it re-synchronizes, how quickly, and with
what characteristic fluctuations --- is the diagnostic signal.

\subsection{Sampling the order parameter}

After each integration step, the order parameter $r$ is evaluated using
equation~\eqref{eq:orderparam} and appended to a trajectory. For a
recording with $M$ beats, the trajectory has approximately
$M \cdot n_{\text{steps}}$ samples.

The CSI reported for the recording is the mean of the trajectory:
\begin{equation}
\mathrm{CSI} = \bar{r} = \frac{1}{T} \sum_{t=1}^{T} r(t),
\label{eq:csi}
\end{equation}
where $T$ is the trajectory length.

\section{Early Warning Signals}\label{sec:ews}

The CSI is a scalar summary of the trajectory $r(t)$. The trajectory
itself carries more information than its mean. Four statistical moments
of $r(t)$ are computed and reported alongside the CSI, as early warning
signals of impending rhythm transitions.

\subsection{Variance}

The sample variance of the trajectory is
\begin{equation}
\mathrm{Var}[r] = \frac{1}{T-1} \sum_{t=1}^{T} (r(t) - \bar{r})^2.
\label{eq:var}
\end{equation}
A rise in the variance of an order parameter is a classic signature of
critical slowing down: as a system approaches a bifurcation, it
recovers more slowly from perturbations and therefore exhibits larger
fluctuations~\cite{scheffer2009}.

\subsection{Skewness}

The sample skewness is
\begin{equation}
\mathrm{Skew}[r] = \frac{T}{(T-1)(T-2)} \sum_{t=1}^{T}
\left( \frac{r(t) - \bar{r}}{\sigma_r} \right)^3,
\label{eq:skew}
\end{equation}
where $\sigma_r = \sqrt{\mathrm{Var}[r]}$. The skewness is a measure of
asymmetry in the distribution of $r(t)$. It reflects the extent to which
the trajectory spends time at values above its mean versus below.

\subsection{Excess kurtosis}

The excess kurtosis is
\begin{equation}
\mathrm{Kurt}[r] = \frac{1}{T} \sum_{t=1}^{T}
\left( \frac{r(t) - \bar{r}}{\sigma_r} \right)^4 - 3.
\label{eq:kurt}
\end{equation}
The excess kurtosis measures the weight of the tails of the
distribution relative to a Gaussian. Positive excess kurtosis indicates
intermittent bursts of incoherence: the trajectory spends most of its
time near its mean but occasionally excursions to extreme values.

\subsection{Lag-1 autocorrelation}

The lag-1 autocorrelation is
\begin{equation}
\rho_1[r] = \frac{\sum_{t=2}^{T} (r(t) - \bar{r})(r(t-1) - \bar{r})}
{\sum_{t=1}^{T} (r(t) - \bar{r})^2}.
\label{eq:lag1}
\end{equation}
The lag-1 autocorrelation approaches 1 as the system recovers more
slowly from perturbations. Together with the variance, it is one of the
most reliable early warning signals identified in the
literature~\cite{scheffer2009}.

\section{Mathematical Properties}\label{sec:properties}

\subsection{Invariance to a common frequency offset}

Adding a constant $c$ to every natural frequency in
equation~\eqref{eq:kuramoto} is equivalent to applying a global rotation
to all phases. To see this, let $\tilde{\theta}_i(t) = \theta_i(t) - ct$.
Then
\begin{equation}
\frac{d\tilde{\theta}_i}{dt} = \frac{d\theta_i}{dt} - c
= (\omega_i - c) + \frac{K}{N}\sum_{j=1}^{N}
\sin(\tilde{\theta}_j - \tilde{\theta}_i),
\label{eq:offset}
\end{equation}
because $\theta_j - \theta_i = \tilde{\theta}_j - \tilde{\theta}_i$.
The transformed phases satisfy the same Kuramoto equation with shifted
frequencies. The order parameter, equation~\eqref{eq:orderparam}, is
invariant to a global rotation, so $r$ is unchanged.

This property is what justifies the frequency normalization of
equation~\eqref{eq:normalization}. Subtracting the mean of the
frequencies does not change the relative dynamics. Scaling by the
standard deviation also does not change the relative dynamics, up to a
time reparameterization: rescaling all frequencies by $s$ is equivalent
to rescaling the time axis by $1/s$.

\subsection{The uniform distribution is a fixed point}

Let all phases be equally spaced around the circle:
\begin{equation}
\theta_i = \frac{2\pi i}{N}, \quad i = 0, 1, \ldots, N-1.
\label{eq:uniform}
\end{equation}
For any $i$, the sum of sines in equation~\eqref{eq:kuramoto} is
\begin{equation}
\sum_{j=0}^{N-1} \sin(\theta_j - \theta_i) = 0,
\label{eq:zero}
\end{equation}
because the pairs $\theta_j$ and $\theta_{j'}$ that are symmetric
around $\theta_i$ contribute opposite sine values. The right-hand side
of equation~\eqref{eq:kuramoto} is therefore just $\omega_i$ for every
$i$.

If all $\omega_i$ are equal, no oscillator moves and the network remains
in the uniform configuration forever, at any coupling strength. In the
thermodynamic limit $N \to \infty$, this configuration is unstable
above $K_c$: any perturbation grows exponentially. But for finite $N$,
and starting exactly at the uniform distribution, the network is a fixed
point.

This is a mathematical property, not a bug. Its practical consequence is
that any test of the coupling mechanism must use an initial condition
that is not perfectly symmetric. In the CSI pipeline, this property is
irrelevant, because the phase resets at each R-peak are inherently
non-symmetric.

\subsection{Dependence on the number of oscillators}

The order parameter is bounded by $r \leq 1$ for any $N$. The expected
value of $r$ for $N$ oscillators with uniformly random phases is
\begin{equation}
\mathbb{E}[r] = \frac{\sqrt{\pi}}{2\sqrt{N}} \cdot
\frac{\Gamma(N)}{\Gamma(N + 1/2)} \approx \frac{1}{\sqrt{N}}
\quad \text{for large } N.
\label{eq:baseline}
\end{equation}
For $N = 16$, the incoherent baseline is approximately $0.25$. CSI
values meaningfully above this baseline indicate partial
synchronization.

\subsection{Residual information beyond RR variability}

The residual of the CSI is defined as the component orthogonal to the
coefficient of variation of the RR series. The companion paper reports
the full analysis; here we note the mathematical result that makes the
residual non-trivial: the CSI is not a linear function of the RR series
coefficient of variation, so the OLS residual is not identically zero.
At $K = 2$, the residual carries information about the atrial state that
the coefficient of variation does not, though the size of that
information depends on the frequency distribution (Section~\ref{sec:network})
and on the coupling strength (Section~\ref{sec:coupling}).

\section{Limitations}\label{sec:limitations}

\subsection{The model is a simplification}

The Kuramoto model describes a population of coupled phase oscillators
with no amplitude, no spatial structure, and sine coupling. The
sinoatrial node is none of these things: it has spatial extent, it has
finite amplitude dynamics, and the coupling between pacemaker cells is
not necessarily sinusoidal. The Kuramoto model is a minimal description,
chosen because it captures the essential phenomenon (phase locking) with
the fewest parameters.

Whether the minimal model is sufficient for the diagnostic purpose is an
empirical question. If the CSI discriminates clinical states and predicts
transitions, the minimality is justified. If it does not, a more detailed
model is needed. Revision 0.2 of the companion residual paper narrows
the claim: the residual carries information beyond standard HRV, but the
information is modest and depends on the frequency distribution.

\subsection{The frequency distribution is arbitrary}

The choice to use the Riemann zeros as the frequency distribution is
inherited from the reference implementation. It is not derived from
cardiac physiology. The sensitivity analysis shows that the trajectory
of the CSI is not robust to this choice at any useful coupling strength,
and that the residual is diagnostic for some distributions but not
others. The frequency distribution is therefore part of the definition
of the CSI, not an incidental parameter.

\subsection{The coupling strength is a free parameter}

$K = 2$ is the only value in the tested range where the residual
separates sinus from AF. The useful range is narrow: below $K = 1$ the
residual inverts in sign, and above $K = 3$ the network saturates. The
choice is empirical, not derived. The sensitivity analysis reports the
full sweep.

\subsection{The single-lead constraint}

The ECG analyzed is single-lead. No spatial information is available.
The CSI is a temporal metric: it describes the evolution of
synchronization over time, not its distribution over the atrium. A
multi-lead device would allow a spatial extension, but that is out of
scope for this application.

\section{Conclusion}

The Kuramoto model provides a mathematically rigorous framework for
describing phase synchronization in coupled systems. Applied to the RR
tachogram of an ECG recording, it yields the Cardiac Synchronization
Index, a phase-domain measure of how regularly the heart is beating.
The model's parameters are documented, its numerical integration is
described, and its limitations are stated. The coupling strength and
the frequency distribution are the two parameters that determine
whether the metric is diagnostic; both are reported in the companion
documents.

The CSI is a research metric. Its clinical validity depends on the
outcomes of the validation protocols described in the companion white
paper and the sensitivity analysis. The mathematics in this document
establishes that the metric is well-defined and computable; it does not
establish that the metric is clinically useful.

\begin{thebibliography}{9}

\bibitem{kuramoto1975}
Y.~Kuramoto. Self-entrainment of a population of coupled non-linear
oscillators. In \emph{International Symposium on Mathematical Problems
in Theoretical Physics}, pages 420--422. Springer, 1975.

\bibitem{acebron2005}
J.~A.~Acebr\'on, L.~L.~Bonilla, C.~J.~P\'erez Vicente, F.~Ritort, and
R.~Spigler. The Kuramoto model: A simple paradigm for synchronization
phenomena. \emph{Reviews of Modern Physics}, 77(1):137--185, 2005.

\bibitem{strogatz2000}
S.~H.~Strogatz. From Kuramoto to Crawford: exploring the onset of
synchronization in populations of coupled oscillators.
\emph{Physica D}, 143(1--4):1--20, 2000.

\bibitem{scheffer2009}
M.~Scheffer, J.~Bascompte, W.~A.~Brock, V.~Brovkin, S.~R.~Carpenter,
V.~Dakos, H.~Held, E.~H.~van Nes, M.~Rietkerk, and G.~Sugihara.
Early-warning signals for critical transitions. \emph{Nature},
461(7260):53--59, 2009.

\bibitem{simoes2026whitepaper}
A.~Simoes. NeuroCardia: An Open-Source, On-Device Phase-Domain
Diagnostic Engine for Single-Lead ECG. White paper, revision 0.2, 2026.
Companion document.

\bibitem{simoes2026sensitivity}
A.~Simoes. Sensitivity Analysis Protocol: Validation Methodology for
the Kuramoto Phase Engine in NeuroCardia. Protocol, revision 0.2, 2026.
Companion document.

\bibitem{simoes2026residual}
A.~Simoes. The Kuramoto Residual: Phase-Domain Information in the
Cardiac Rhythm Beyond Heart Rate Variability. Research paper, revision
0.2, 2026. Companion document.

\end{thebibliography}

\end{document}
